\chapter{Two Ways to Grow: Shape or Value}
\label{ch:erweiterungen}

% Register: D5, S8. Sources: spec. sections 5 and 6.

\section{The moment of the fork}

Take the seven in its home $H_4$: there it is written $0012$. Now
we give it one more place. What is that supposed to mean?

The first answer is the naive one: we put a zero in front.
\[
0012 \;\longmapsto\; 00012 .
\]
The digits stay put -- but in $H_5$ every place except the last
carries a different weight, and the value slips from $7$ to $8$.
The \emph{shape} was preserved, the value was not.

The second answer is the cautious one: we want the number to stay
the same, and we re-encode the $7$ in $H_5$:
\[
0012 \;\longmapsto\; 00011 .
\]
The \emph{value} was preserved, the shape was not.

In the decimal system both answers are the same operation -- which
is why we never noticed that they are two questions. In Hori space
they part ways, and we give them names: the structure-preserving
-- equivalently we say: shape-preserving -- extension $Z$ puts the
zero in front, the value-preserving lift $L$ re-encodes.

\begin{definition}[Register D5]
$Z_n(a_1,\ldots,a_n)=(0,a_1,\ldots,a_n)$ and
$L_n(a)=E_{n+1}(V_n(a))$, where $E_{n+1}$ is the encoding into
$\Hn{n+1}$. We have $V_{n+1}(L_n(a))=V_n(a)$, whereas $Z_n$ in
general changes the value.
\end{definition}

\section{Two staircases, two skies}

Both operations join the finite horizons into an infinite
staircase -- but they are \emph{different} staircases, and they
lead to different places. The question in each case is: which
elements of consecutive horizons count as ``the same''?

Under $L$ two representations count as the same if they have the
same value: $013$ in $H_3$, $0012$ in $H_4$, $00011$ in $H_5$ are
a single object -- the seven. Under $Z$ they count as the same if
they have the same shape: $12$, $012$, $0012$ are a single object
-- the pattern ``one-two'', whose value wanders from step to step.

\begin{satz}[Register S8]\label{satz:grenzen}
The limit space of the staircase under $L$ is $\Nn$, the set of
ordinary natural numbers. The limit space under $Z$ is the space
of all shapes -- formally: the semiring of polynomials with
non-negative integer coefficients in the basis of falling
factorials, which the next chapter introduces.
\end{satz}

\begin{proof}[Proof idea]
Under $L$ precisely the representations of the same value are
identified (the encoding is injective), and every value occurs --
the classes \emph{are} the natural numbers. Under $Z$ precisely
those digit strings are identified that differ by leading zeros,
and every shape occurs from its smallest admissible horizon
onwards -- the classes are the shapes. The complete proof via
directed systems is given in Appendix~\ref{ch:anhang-beweise}.
\end{proof}

\begin{beispiel}
The two staircases, side by side. Preserving value, the seven
remains itself and changes shape; preserving shape, the pattern
``one-two'' remains and changes value:
\begin{center}
\begin{tabular}{l|llll}
World & $H_2$ & $H_3$ & $H_4$ & $H_5$\\
\hline
value-preserving (always 7) & --- & $013$ & $0012$ & $00011$\\
shape-preserving (always ``12'') & $12=5$ & $012=6$ & $0012=7$ & $00012=8$\\
\end{tabular}
\end{center}
Note the crossing: $0012$ appears in both rows -- once as the
seven's second home, once as the pattern's third life. The same
element, two readings.
\end{beispiel}

This is the first theorem in this book that one may read
philosophically: \emph{the same finite spaces generate different
infinities, depending on which notion of consistency one chooses
between the steps.} The familiar number line is one of these
infinities -- the value-preserving one. The other, the
shape-preserving one, is mathematically just as legitimate; only,
in the ordinary number system nobody has ever been able to tell it
apart from the first, because there the two staircases coincide.

\section{The third staircase}

For a long time we thought there were exactly two ways to grow.
Then an outside mathematician read a short description of this
system -- and inadvertently anchored the ladder of bases at the
other end. Out of his ``misunderstanding'' fell a third staircase.

For one can also append a zero to a number \emph{on the right}.
All digits then keep their position counted from the left, and a
strikingly clean rule holds: the value does not multiply out of
control but scales \emph{exactly} by the factor $n+2$ --
$12$ in $H_2$ (value $5$) becomes $120$ in $H_3$ (value $20$).
What stays invariant is neither the value nor the shape but a
third quantity: the \emph{factorial fraction}
$F=\sum_i a_i/(i+1)!$, since we always have
\[
V_n(a)=(n+1)!\cdot F(a).
\]
Every Hori number is thus also a fraction between zero and one,
times the size of its world -- and this fraction system is an old
acquaintance: it is Cantor's factorial base for fractions, a good
century and a half old.

With that, Hori space has three staircases, three invariants,
three skies:

\begin{center}
\begin{tabular}{l|lll}
Extension & Invariant & Value & Limit space\\
\hline
value-preserving ($L$) & value & $\times1$ & natural numbers\\
shape-preserving ($Z$, left) & polynomial & uneven & structure space\\
fraction-preserving ($R$, right) & fraction $F$ & $\times(n+2)$ & Cantor fractions\\
\end{tabular}
\end{center}

And the objection ``but this has been known since Cantor'' thereby
receives a precise answer: it hits the third staircase completely
-- and the other two, together with everything this book proves
about their interplay, not at all.

\section{The eigenhorizon: where the fork heals}

For every single number there is a point at which the difference
between $Z$ and $L$ vanishes. If the entire value sits in the last
digit -- the number $n$ as $0\cdots0n$; we call this shape its
\emph{terminal representation}, and it becomes possible from the
horizon $H_n$ onwards --, then the zero extension is
value-neutral: shape and value grow together from there on. We
call $H_n$ the \emph{eigenhorizon} of $n$. Below it the number is
rebuilt at every move; from there on it only grows quietly to the
left. One can compute exactly when a leading zero is
value-neutral: precisely when all digits except the last are zero
-- and the next chapter turns this into a one-liner using the
horizon derivative.

What this feels like for a concrete number is told in the
interlude ``The life story of the number seven''.
